\SourceLeaf{63}
\begin{center}\emph{A method for finding primitive roots.}\end{center}

\Article{73}
Methods for finding primitive roots rely largely on trial. Comparing what we gave in Article 55 with what we will give below on solving the congruence $x^n\equiv1$ provides almost everything that direct methods can achieve. Euler admits, in \emph{Opusc. Analyt. T.} I, \emph{p.} 152, that finding these numbers seems extremely difficult, and that their nature must be counted among the deepest mysteries of numbers. But they can be determined quite readily by trial, as follows. With experience, one learns to shorten the lengthy work by various special devices; these are learned much faster through practice than through rules.

$1^0$. Take any number coprime to $p$ (we always use this letter for the modulus), say $a$. It usually shortens the calculation to choose as small a number as possible, \emph{for example} $2$. Determine its period (Article 46), \emph{that is,} the least residues of its powers, until a power $a^t$ is reached whose least residue is $1$\HistoricalNote{*)}{One readily sees that the powers themselves need not be known, since the least residue of each can easily be obtained from the least residue of the preceding power.}. If $t=p-1$, then $a$ is a primitive root.

$2^0$. But if $t<p-1$, take another number $b$ not contained in the period of $a$, and find its period in the same way. Denote the exponent to which $b$ belongs by $u$. It is easy to see that $u$ can neither equal $t$ nor be a proper divisor of it: in either case we would have $b^t\equiv1$, which is impossible, since the period of $a$ includes all numbers whose power with exponent $t$ is congruent to unity (Article 53). If $u$ is $=p-1$, then $b$ is a primitive root. If $u$ is not $=p-1$ but is a multiple of $t$, we have gained a number belonging to a larger exponent, and have moved closer to our aim of finding a number belonging to the \emph{largest} exponent. If $u$ is neither $=p-1$ nor a multiple of $t$, we can still find a number belonging to an exponent larger than $t,u$, namely the least common multiple of $t,u$. Let this be $=y$, and factor $y$ into two coprime factors $m,n$, one dividing $t$ and the other $u$\HistoricalNote{\ensuremath{\dagger})}{Article 18 readily shows how this can be done. Factor $y$ into distinct primes or powers of distinct primes. Each factor divides one of the numbers $t,u$, or both. Assign each to $t$ or $u$, according to which it divides; if it divides both, either assignment may be made. Let the product of those assigned to $t$ be $=m$, and the product of the rest $=n$. It is then easy to see that $m$ divides $t$, $n$ divides $u$, and $mn=y$.}. Now take the power with exponent $\frac{t}{m}$ of $a$, putting it $\equiv A$, and the
\SourceLeaf{64}power with exponent $\frac{u}{n}$ of $b$, putting it $\equiv B\,(\operatorname{mod.}p)$. The product $AB$ will be a number belonging to the exponent $n$. For $A$ clearly belongs to the exponent $m$, and $B$ to $n$. So the product $AB$ belongs to $mn$, since $m,n$ are coprime. This can be proved in exactly the same way as in Article 55, II.

$3^0$. If now $y=p-1$, then $AB$ is a primitive root. Otherwise, as before, another number not in the period of $AB$ must be used. This is either a primitive root, or belongs to an exponent larger than $y$, or at least allows us, as before, to find a number belonging to an exponent larger than $y$. Since the numbers obtained by repeating this operation belong to continually increasing exponents, we must eventually find one belonging to the \emph{largest} exponent, \emph{that is,} a primitive root, \emph{as required.}

\Needspace{7\baselineskip}
\Article{74}
An example will make these rules clearer. Let $p=73$, and seek a primitive root for it. First try $2$, whose period is

\[\setlength{\arraycolsep}{2.5pt}\begin{array}{@{}rrrrrrrrrrl@{}}
1.&2.&4.&8.&16.&32.&64.&55.&37.&1&\text{etc.}\\
0.&1.&2.&3.&4.&5.&6.&7.&8.&9&\text{etc.}
\end{array}\]

Since the power with exponent $9$ is already congruent to unity, $2$ is not a primitive root. Try another number not in the period of $2$, \emph{for example} $3$, whose period is

\[\setlength{\arraycolsep}{2.5pt}\begin{array}{@{}rrrrrrrrrrrrrl@{}}
1.&3.&9.&27.&8.&24.&72.&70.&64.&46.&65.&49.&1&\text{etc.}\\
0.&1.&2.&3.&4.&5.&6.&7.&8.&9.&10.&11.&12&\text{etc.}
\end{array}\]

So $3$ is not a primitive root either. But the least common multiple of the exponents to which $2,3$ belong, \emph{that is,} of $9,12$, is $36$. By the rules of the preceding article, this factors into $9$ and $4$. We must therefore raise $2$ to the power with exponent $\frac{9}{9}$, \emph{that is,} keep $2$ itself, and raise $3$ to the power with exponent $3$. Their product is $54$, which therefore belongs to the exponent $36$. Finally, calculating the period of $54$ and again trying a number not in it, \emph{for example} $5$, will show that this number is a primitive root.

\SourceLeaf{65}
\begin{center}\emph{Various theorems on periods and primitive roots.}\end{center}

\Article{75}
Before leaving this subject, we will give a few propositions whose simplicity seems to make them worth attention.

\emph{The product of all the terms in the period of any number is $\equiv1$ when their number, or the exponent to which the number belongs, is odd, and $\equiv-1$ when that exponent is even.}

\emph{Example.} For the modulus $13$, the period of $5$ consists of $1,5,12,8$, whose product is $480\equiv-1\,(\operatorname{mod.}13)$.

For the same modulus, the period of $3$ consists of $1,3,9$, whose product is $27\equiv1\,(\operatorname{mod.}13)$.

\emph{Proof.} Let the exponent to which the number belongs be $t$, and its index $\frac{p-1}{t}$. This is always possible with a suitable choice of base (Article 71). Then the index of the product of all terms of the period is

\[\equiv(1+2+3+\text{etc.}+t-1)\frac{p-1}{t}=\frac{(t-1)(p-1)}{2}\]

\emph{that is,} $\equiv0\,(\operatorname{mod.}p-1)$ when $t$ is odd, and $\equiv\frac{p-1}{2}$ when $t$ is even. Thus in the first case the product $\equiv1\,(\operatorname{mod.}p)$, and in the second $\equiv-1\,(\operatorname{mod.}p)$ (Article 62). \mbox{\emph{Q. E. D.}}

\Article{76}
If the number in the preceding theorem is a primitive root, its period includes all the numbers $1,2,3\mathpunct{.\,.\,.\,.\,.\,.}\,p-1$, whose product is therefore always $\equiv-1$ (for $p-1$ is always even, apart from the single case $p=2$, where $-1$ and $+1$ are equivalent). This elegant theorem, usually stated as follows: \emph{the product of all numbers smaller than a given prime, increased by unity, is divisible by that prime}, was first stated by Waring and attributed to Wilson, Esq., \emph{Meditt. algebr. Ed.} 3, \emph{p.} 380. But neither could prove it. Waring admits that the proof seems all the harder because no \emph{notation} can be devised to express a prime number. --- In our view, however, truths of this kind should be drawn from concepts rather than notations. Later La Grange gave a proof, \emph{Nouv. M\'em. de l'Ac. de Berlin,} 1771. It rests on considering
\SourceLeaf{66}the coefficients obtained by expanding the product

\[x+1.\,x+2.\,x+3\mathpunct{.\,.\,.\,.}\,x+p-1\]

\noindent Namely, if this product is put

\[\equiv x^{p-1}+Ax^{p-2}+Bx^{p-3}+\text{etc.}+Mx+N\]

\noindent the coefficients $A,B$, etc., $M$ will be divisible by $p$, while $N$ will be $=1.2.3\mathpunct{.\,.\,.\,.}\,p-1$. Now for $x=1$ the product is divisible by $p$, but is then $\equiv1+N\,(\operatorname{mod.}p)$. So $1+N$ must be divisible by $p$.

Finally, Euler gave a proof in \emph{Opusc. analyt. T.} I, \emph{p.} 329, agreeing with the one we have given here. If such men thought this theorem worth their consideration, we hope readers will not mind if we add another proof.

\Article{77}
When modulo $p$ the product of two numbers $a,b$ is congruent to unity, let us follow Euler in calling the numbers $a,b$ \emph{partners}. By the preceding section, every positive number smaller than $p$ has exactly one positive partner smaller than $p$. It is easy to prove that among $1,2,3\mathpunct{.\,.\,.\,.}\,p-1$, only $1$ and $p-1$ are their own partners. For numbers that are their own partners are roots of $xx\equiv1$, which, being of the second degree, cannot have more than two roots, \emph{that is,} any besides $1$ and $p-1$. Removing these, the remaining numbers $2.3\mathpunct{.\,.\,.\,.}\,p-2$ are paired off. Their product is therefore $\equiv1$, so the product of all the numbers $1,2,3\mathpunct{.\,.\,.}\,p-1$ is $\equiv p-1$, or $\equiv-1$. \mbox{\emph{Q. E. D.}}

\emph{For example,} for $p=13$ the numbers $2,3,4\mathpunct{.\,.\,.\,.}\,11$ are paired as follows: $2$ with $7$; $3$ with $9$; $4$ with $10$; $5$ with $8$; $6$ with $11$. Namely, $2.7\equiv1$; $3.9\equiv1$, etc. So $2.3.4\mathpunct{.\,.\,.\,.}\,11\equiv1$, and therefore $1.2.3\mathpunct{.\,.\,.\,.}\,12\equiv-1$.

\Needspace{9\baselineskip}
\Article{78}
Wilson’s theorem can be stated more generally as follows. \emph{The product of all numbers smaller than any given number $A$ and coprime to it is congruent modulo $A$ to unity, taken either negatively or positively.} Unity is to be taken negatively when $A$ has the form $p^m$ or $2p^m$, where $p$ denotes a
\SourceLeaf{67}prime other than $2$, and also when $A=4$; positively in every other case. The theorem as stated by Wilson is included in the first case. --- \emph{For example,} for $A=15$ the product of $1,2,4,7,8,11,13,14$ is $\equiv1\,(\operatorname{mod.}15)$. For brevity we give no proof. We only note that it can be carried out as in the preceding article, except that $xx\equiv1$ may have more than two roots, requiring some special considerations. A proof could also be obtained by considering indices, as in Article 75, together with what we will soon give on nonprime moduli.

\Needspace{8\baselineskip}
\Article{79}
We return to listing other propositions (Article 75).

\emph{The sum of all terms in the period of any number is $\equiv0$}, as in the example of Article 75, $1+5+12+8=26\equiv0\,(\operatorname{mod.}13)$.

\emph{Proof.} Let the number whose period is in question be $=a$, and the exponent to which it belongs $=t$. The sum of all terms of the period is then
\[\equiv1+a+aa+a^3+\text{etc.}+a^{t-1}\equiv\frac{a^t-1}{a-1}\,(\operatorname{mod.}p)\]
But $a^t-1\equiv0$, so this sum is always $\equiv0$ (Article 22), unless $a-1$ is divisible by $p$, or $a\equiv1$. We must therefore exclude this case if we wish to call even a single term a \emph{period}.

\Article{80}
\emph{The product of all primitive roots is $\equiv1$}, with the single exception $p=3$, when there is just one primitive root, $2$.

\emph{Proof.} With any primitive root taken as the base, the indices of all primitive roots are the numbers coprime to $p-1$ and smaller than it. The sum of these numbers, \emph{that is,} the index of the product of all primitive roots, is $\equiv0\,(\operatorname{mod.}p-1)$, so the product $\equiv1\,(\operatorname{mod.}p)$. For if $k$ is coprime to $p-1$, clearly $p-1-k$ is also coprime to $p-1$. Thus the numbers coprime to $p-1$ pair off into sums divisible by $p-1$. (But $k$ can never equal $p-1-k$, apart from the case $p-1=2$, or $p=3$, which we have excluded: in every other case $\frac{p-1}{2}$ is clearly not coprime to $p-1$.)

\SourceLeaf{68}
\Article{81}
\emph{The sum of all primitive roots is either} $\equiv0$ (when $p-1$ is divisible by a \emph{square}) \emph{or $\equiv\pm1\,(\mathit{mod.}p)$} (when $p-1$ is a product of distinct primes; take the positive sign if their number is even and the negative sign if it is odd).

\emph{Examples.} $1^0$ For $p=13$, the primitive roots are $2,6,7,11$, whose sum is $26\equiv0\,(\operatorname{mod.}13)$.

$2^0$ For $p=11$, the primitive roots are $2,6,7,8$, whose sum is $23\equiv+1\,(\operatorname{mod.}11)$.

$3^0$ For $p=31$, the primitive roots are $3,11,12,13,17.\,21,22,24$, whose sum is $123\equiv-1\,(\operatorname{mod.}31)$.

\emph{Proof.} We showed above (Article 55, II) that if $p-1$ is $=a^{\alpha}b^{\beta}c^{\gamma}$, etc. (where $a,b,c$, etc., denote distinct primes), and $A,B,C$, etc., are any numbers belonging respectively to the exponents $a^{\alpha},b^{\beta},c^{\gamma}$, etc., all products $ABC$, etc., are primitive roots. It is also easy to prove that every primitive root can be represented by such a product, and in only one way\HistoricalNote{*)}{Determine the numbers $\mathfrak{a},\mathfrak{b},\mathfrak{c}$, etc., so that $\mathfrak{a}\equiv1\,(\operatorname{mod.}a^{\alpha})$ and $\equiv0\,(\operatorname{mod.}b^{\beta}c^{\gamma}\text{etc.})$; $\mathfrak{b}\equiv1\,(\operatorname{mod.}b^{\beta})$ and $\equiv0\,(\operatorname{mod.}a^{\alpha}c^{\gamma}\text{etc.})$, etc. (see Article 32). Then $\mathfrak{a}+\mathfrak{b}+\mathfrak{c}+\text{etc.}\equiv1\,(\operatorname{mod.}p-1)$ (Article 19). Now if any primitive root $r$ is to be represented by a product $ABC$, etc., take $A\equiv r^{\mathfrak{a}},B\equiv r^{\mathfrak{b}},C\equiv r^{\mathfrak{c}}$, etc. Then $A$ belongs to the exponent $a^{\alpha}$, $B$ to the exponent $b^{\beta}$, etc. The product of all $A,B,C$, etc., is $\equiv r\,(\operatorname{mod.}p)$. Finally, it is easy to see that $A,B,C$, etc., cannot be determined in any other way.}.

So these products can take the place of the primitive roots themselves. Since in these products every value of $A$ must be combined with every value of $B$, etc., the sum of all the products equals the product of the sum of all values of $A$, the sum of all values of $B$, the sum of all values of $C$, etc., as is known from the theory of combinations. Denote all values of $A$; $B$, etc., by $A,A',A''$, etc.; $B,B',B''$, etc., etc. Then the sum of all primitive roots
\[\equiv(A+A'+\text{etc.})(B+B'+\text{etc.})\text{etc.}\]
Now I say that if the exponent $\alpha$ is $=1$, the sum $A+A'+A''+\text{etc.}$ will be $\equiv-1\,(\operatorname{mod.}p)$; but if $\alpha$ is $>1$, this sum will be $\equiv0$. The same applies to the other exponents $\beta,\gamma$, etc. Once this is proved, the truth of our theorem will
\SourceLeaf{69}be evident. For when $p-1$ is divisible by a square, one of the exponents $\alpha,\beta,\gamma$, etc., exceeds unity. So one of the factors whose product is congruent to the sum of all primitive roots is $\equiv0$, and therefore so is the product itself. But when $p-1$ is divisible by no square, all the exponents $\alpha,\beta,\gamma$, etc., are $=1$. The sum of all primitive roots is then congruent to a product with as many factors, each $\equiv-1$, as there are numbers $a,b,c$, etc. It is therefore $\equiv\pm1$, according to whether their number is even or odd. The claims are proved as follows.

$1^0$. When $\alpha=1$ and $A$ is a number belonging to the exponent $a$, the other numbers belonging to this exponent are $A^2,A^3\GaussFourDots A^{a-1}$. But
\[1+A+A^2+A^3\GaussFourDots+A^{a-1}\]
is the sum of a complete period, and therefore $\equiv0$ (Article 79). So
\[A+A^2+A^3\GaussFourDots+A^{a-1}\equiv-1\]

$2^0$. But when $\alpha>1$ and $A$ is a number belonging to the exponent $a^{\alpha}$, the other numbers belonging to this exponent are obtained by removing from $A^2,A^3,A^4\GaussFourDots A^{a^{\alpha}-1}$ the terms $A^a,A^{2a},A^{3a}$, etc.; see Article 53. Their sum is therefore
\[\equiv1+A+A^2\GaussFourDots+A^{a^{\alpha}-1}-(1+A^a+A^{2a}\GaussFourDots+A^{a^{\alpha}-a})\]
\emph{that is,} congruent to the difference of two periods, and thus $\equiv0$. \mbox{\emph{Q. E. D.}}

\begin{center}\emph{Moduli that are powers of primes.}\end{center}

\Article{82}
Everything given so far rests on the assumption that the modulus is prime. We still have to consider the case where a composite number is taken as the modulus. But the properties here are neither as elegant as in the previous case, nor do they require subtle devices to find them. Rather, almost everything can be drawn simply from applying the preceding principles. Going through every detail would therefore be unnecessary and tedious. So we will briefly give what this case shares with the previous one and what is particular to it.

\SourceLeaf{70}
\Needspace{10\baselineskip}
\Article{83}
The propositions of Articles 45–48 have already been proved in general. But the proposition of Article 49 must be modified as follows:

\emph{If $f$ denotes how many numbers are coprime to $m$ and smaller than $m$, that is, if $f=\phi m$ (Article 38), the exponent $t$ of the lowest power of a given number $a$ coprime to $m$ that is congruent to unity modulo $m$ is either $=f$ or a proper divisor of that number.}

The proof in Article 49 also applies here if we replace $p$ everywhere by $m$, $p-1$ by $f$, and the numbers $1,2,3,\GaussFourDots p-1$ by the numbers coprime to $m$ and smaller than $m$. We therefore refer the reader there. The other proofs mentioned there (Articles 50, 51) cannot be applied to this case without considerable complications. --- But for the following propositions, Article 52 onward, a major difference appears between moduli that are powers of primes and those divisible by several primes. So we will consider moduli of the first kind separately.

\Article{84}
If the modulus $m=p^n$, where $p$ is prime, then $f=p^{n-1}(p-1)$ (Article 38). Applying the investigations of Articles 53, 54 to this case, with the changes prescribed in the preceding article, shows that everything proved there also holds here, provided we have first proved that a congruence of the form $x^t-1\equiv0\,(\operatorname{mod.}p^n)$ cannot have more than $t$ different roots. For a prime modulus we derived this from the more general proposition of Article 43. But that proposition, in its full extent, holds only for prime moduli, and so cannot be applied here. Nevertheless, we will prove the proposition for this particular case by a special method. Below, in Section VIII, we will show how to reach the same result more easily.

\Article{85}
We set out to prove the following theorem:

\emph{If the greatest common divisor of $t$ and $p^{n-1}(p-1)$ is $e$, the congruence $x^t\equiv1\,(\mathit{mod.}p^n)$ has $e$ different roots.}

Let $e=kp^{\nu}$, where $k$ does not contain the factor $p$, and therefore divides $p-1$.
\SourceLeaf{71}Then the congruence $x^t\equiv1$ has, modulo $p$, exactly $k$ different roots. Denote these by $A,B,C$, etc. Any root of the same congruence modulo $p^n$ must be congruent modulo $p$ to one of $A,B,C$, etc. We will now show that $x^t\equiv1.\,(\operatorname{mod.}p^n)$ has $p^{\nu}$ roots congruent to $A$, as many congruent to $B$, etc., modulo $p$. The total number of roots will then be $kp^{\nu}$, or $e$, as asserted. We will arrange the proof as follows. \emph{First,} we show that if $\alpha$ is a root congruent to $A$ modulo $p$, then

\[\alpha+p^{n-\nu},\ \alpha+2p^{n-\nu},\ \alpha+3p^{n-\nu},\mathinner{.\,.\,.\,.}\alpha+(p^{\nu}-1)p^{n-\nu}\]

are also roots. \emph{Second,} among numbers congruent to $A$ modulo $p$, no numbers besides those of the form $\alpha+hp^{n-\nu}$ (where $h$ is any integer) can be roots. This clearly gives $p^{\nu}$ different roots and no more. The same holds for roots congruent to each of $B,C$, etc. \emph{Third,} we show how a root congruent to $A$ modulo $p$ can always be found.

\Needspace{10\baselineskip}
\Article{86}
\textsc{Theorem.} \emph{If, as in the preceding article, $t$ is divisible by $p^{\nu}$ but not by $p^{\nu+1}$, then}

\[\begin{aligned}
(\alpha+hp^{\mu})^t-\alpha^t&\equiv0\,(mod.p^{\mu+\nu}),\\
&\textit{but}\ \equiv\alpha^{t-1}hp^{\mu}t\,(mod.p^{\mu+\nu+1})
\end{aligned}\]

The latter part of the theorem does not hold when both $p=2$ and $\mu=1$.

We could prove this theorem by expanding the power of the binomial and showing that every term after the second is divisible by $p^{\mu+\nu+1}$. But examining the denominators of the coefficients leads to some complications, so we prefer the following method.

\emph{First} suppose $\mu>1$ and $\nu=1$. Because

\[x^t-y^t=(x-y)(x^{t-1}+x^{t-2}y+x^{t-3}y^2+\text{etc.}+y^{t-1})\]

we have

\[\begin{aligned}
(\alpha+hp^{\mu})^t-\alpha^t={}&hp^{\mu}((\alpha+hp^{\mu})^{t-1}\\
&+(\alpha+hp^{\mu})^{t-2}\alpha+\text{etc.}+\alpha^{t-1})
\end{aligned}\]

But $\alpha+hp^{\mu}\equiv\alpha\,(\operatorname{mod.}p^2)$,
\SourceLeaf{72}so each term $(\alpha+hp^{\mu})^{t-1}$, $(\alpha+hp^{\mu})^{t-2}\alpha$, etc., is $\equiv\alpha^{t-1}(\operatorname{mod.}p^2)$. Thus their sum is $\equiv t\alpha^{t-1}(\operatorname{mod.}p^2)$, or of the form $t\alpha^{t-1}+Vp^2$, where $V$ is some number. Hence $(\alpha+hp^{\mu})^t-\alpha^t$ has the form

\[\begin{aligned}
&\alpha^{t-1}hp^{\mu}t+Vhp^{\mu+2},\\
&\textit{that is,}\ \equiv\alpha^{t-1}hp^{\mu}t(\operatorname{mod.}p^{\mu+2})\quad\text{and}\ \equiv0\,(\operatorname{mod.}p^{\mu+1}).
\end{aligned}\]

\noindent So the theorem is proved in this case.

Now if the theorem failed for other values of $\nu$, still with $\mu>1$, there would necessarily be a boundary up to which it was always true and beyond which it was false. Let the smallest value of $\nu$ for which it is false be $=\varphi$. We readily see that if $t$ is divisible by $p^{\varphi-1}$ but not by $p^{\varphi}$, the theorem is still true, but becomes false when $t$ is replaced by $tp$. So we have

\[\begin{aligned}
(\alpha+hp^{\mu})^t&\equiv\alpha^t+\alpha^{t-1}hp^{\mu}t(\operatorname{mod.}p^{\mu+\varphi})\\
&\text{or}\ =\alpha^t+\alpha^{t-1}hp^{\mu}t+up^{\mu+\varphi}
\end{aligned}\]

\noindent where $u$ is an integer. But since the theorem is already proved for $\nu=1$, we have

\[\begin{aligned}
(\alpha^t+\alpha^{t-1}hp^{\mu}t+up^{\mu+\varphi})^p\equiv{}&\alpha^{tp}+\alpha^{tp-1}hp^{\mu+1}t\\
&+\alpha^{tp-t}up^{\mu+\varphi+1}\,(\operatorname{mod.}p^{\mu+\varphi+1})
\end{aligned}\]

\noindent and therefore also

\[(\alpha+hp^{\mu})^{tp}\equiv\alpha^{tp}+\alpha^{tp-1}hp^{\mu}tp\,(\operatorname{mod.}p^{\mu+\varphi+1})\]

\noindent \emph{that is,} the theorem is true even when $t$ is replaced by $tp$, \emph{that is,} for $\nu=\varphi$, contrary to the hypothesis. Thus the theorem is true for all values of $\nu$.

\Needspace{10\baselineskip}
\Article{87}
The case $\mu=1$ remains. By a method entirely like the one used in the preceding article, without the binomial theorem, we can show that

\[\begin{aligned}
(\alpha+hp)^{t-1}&\equiv\alpha^{t-1}+\alpha^{t-2}(t-1)hp\quad(\operatorname{mod.}p^2)\\
\alpha(\alpha+hp)^{t-2}&\equiv\alpha^{t-1}+\alpha^{t-2}(t-2)hp\\
\alpha\alpha(\alpha+hp)^{t-3}&\equiv\alpha^{t-1}+\alpha^{t-2}(t-3)hp\\
\text{etc.}\quad&
\end{aligned}\]

\noindent so their sum, since the number of terms $=t$, is

\SourceLeaf{73}
\[\equiv t\alpha^{t-1}+\tfrac{(t-1)t}{2}\alpha^{t-2}hp(\operatorname{mod.}p^2)\]

Since $t$ is divisible by $p$, $\dfrac{(t-1)t}{2}$ is also divisible by $p$ in every case except $p=2$, as already noted in the preceding article. In the other cases $\tfrac{(t-1)t}{2}\alpha^{t-2}hp\equiv0\,(\operatorname{mod.}p^2)$, so that sum is also $\equiv t\alpha^{t-1}(\operatorname{mod.}p^2)$, as in the preceding article. The rest of the proof proceeds exactly as there.

We therefore conclude in general, with the single exception $p=2$, that

\[(\alpha+hp^{\mu})^t\equiv\alpha^t(\operatorname{mod.}p^{\mu+\nu})\]

\noindent and $(\alpha+hp^{\mu})^t\ \textit{not}\equiv\alpha^t$ for any modulus that is a higher power of $p$ than $p^{\mu+\nu}$, provided $h$ is not divisible by $p$ and $p^{\nu}$ is the highest power of $p$ dividing $t$.

The first and second propositions we set out to prove in Article 85 follow at once:

\emph{first,} if $\alpha^t\equiv1$, then also $(\alpha+hp^{n-\nu})^t\equiv1\,(\operatorname{mod.}p^n)$;

\emph{second,} suppose a number $\alpha'$, congruent to $A$ and therefore also to $\alpha$ modulo $p$, but not congruent to the latter modulo $p^{n-\nu}$, satisfied $x^t\equiv1\,(\operatorname{mod.}p^n)$. Put $\alpha'$ equal to $=\alpha+lp^{\lambda}$, with $l$ not divisible by $p$. Then $\lambda<n-\nu$, and $(\alpha+lp^{\lambda})^t$ is congruent modulo $p^{\lambda+\nu}$ to $\alpha^t$, but not modulo $p^n$, which is a higher power. So $\alpha'$ cannot be a root of $x^t\equiv1$.

\Article{88}
The \emph{third} task was to find a root of $x^t\equiv1\,(\operatorname{mod.}p^n)$ congruent to $A$. Here we show only how this can be done if a root of the same congruence modulo $p^{n-1}$ is already known. This clearly suffices: starting with the modulus $p$, for which $A$ is a root, we can pass to the modulus $p^2$ and then to every successive power.

So let $\alpha$ be a root of $x^t\equiv1\,(\operatorname{mod.}p^{n-1})$, and seek a root of the same congruence modulo $p^n$. Put this root $=\alpha+hp^{n-\nu-1}$; the preceding article shows that it must have this form. (The case $\nu=n-1$
\SourceLeaf{74}will be considered separately later; $n-1$ cannot be exceeded by $\nu$.) We must therefore have

\[(\alpha+hp^{n-\nu-1})^t\equiv1(\operatorname{mod.}p^{n-1})\]

\noindent But\hfill $(\alpha+hp^{n-\nu-1})^t\equiv\alpha^t+\alpha^{t-1}htp^{n-\nu-1}\,(\operatorname{mod.}p^n)$\hfill\mbox{}

\noindent So if $h$ is chosen so that $1\equiv\alpha^t+\alpha^{t-1}htp^{n-\nu-1}\,(\operatorname{mod.}p^n)$, or, since by hypothesis $1\equiv\alpha^t(\operatorname{mod.}p^{n-1})$ and $t$ is divisible by $p^{\nu}$, so that $\dfrac{\alpha^t-1}{p^{n-1}}+\alpha^{t-1}h\dfrac{t}{p^{\nu}}$ is divisible by $p$, the requirement is met. The preceding section shows that this is always possible, since we assume here that $t$ cannot be divided by a power of $p$ higher than $p^{\nu}$, and so $\alpha^{t-1}\dfrac{t}{p^{\nu}}$ is coprime to $p$.

But if $\nu=n-1$, \emph{that is,} $t$ is divisible by $p^{n-1}$ or even by a higher power of $p$, every value $A$ satisfying $x^t\equiv1$ modulo $p$ also satisfies it modulo $p^n$. For let $t=p^{n-1}\tau$; then $t\equiv\tau(\operatorname{mod.}p-1)$. Since $A^t\equiv1\,(\operatorname{mod.}p)$, we therefore also have $A^{\tau}\equiv1\,(\operatorname{mod.}p)$. Put $A^{\tau}=1+hp$; then $A^t=(1+hp)^{p^{n-1}}\equiv1(\operatorname{mod.}p^n)$ by Article 87.

\Article{89}
Everything derived in Article 57 onward with the help of the theorem that $x^t\equiv1$ cannot have more than $t$ different roots also holds for a modulus that is a power of a prime. If we call the numbers belonging to the exponent $p^{n-1}(p-1)$ \emph{primitive roots}, or equivalently the numbers whose periods contain all numbers not divisible by $p$, primitive roots exist here too. Everything given above on indices and their use, and on solving $x^t\equiv1$, can also be applied here. Since there are no difficulties in doing so, repeating everything from the beginning would be unnecessary. We have also shown how to derive the roots of $x^t\equiv1$ modulo $p^n$ from those of the same congruence modulo $p$. But we must add a few things about moduli that are powers of 2, since they were excluded above.

\Needspace{12\baselineskip}
\begin{center}\emph{Moduli that are powers of two.}\end{center}

\Article{90}
\emph{If a power of 2 higher than the second, say $2^n$, is taken as the modulus, the power with exponent $2^{n-2}$ of any odd number is congruent to unity.}
\SourceLeaf{75}\emph{For example,} $3^8=6561\equiv1\,(\operatorname{mod.}32)$.

Every odd number has one of the forms $1+4h$ or $-1+4h$, so the proposition follows at once from the theorem of Article 86.

Since the exponent to which any odd number belongs modulo $2^n$ must divide $2^{n-2}$, it belongs to one of $1,2,4,8\mathinner{.\,.\,.\,.}2^{2n-2}$. Which one can be decided easily as follows. Let the given number be $=4h\pm1$, and the exponent of the highest power of 2 dividing $h$ be $=m$ (which may also be $=0$, namely when $h$ is odd). Then the exponent to which the given number belongs is $=2^{n-m-2}$, provided $n>m+2$. But if $n=$ or $<m+2$, the given number is $\equiv\pm1$, and so belongs either to the exponent 1 or to the exponent 2. For a number of the form $\pm1+2^{m+2}k$ (equivalent to $4h\pm1$), raised to the power with exponent $2^{n-m-2}$, is congruent to unity modulo $2^n$, but is not when raised to a power whose exponent is a lower power of 2. This follows readily from Article 86. So any number of the form $8k+3$ or $8k+5$ belongs to the exponent $2^{n-2}$.

\Article{91}
Thus there are no \emph{primitive roots} here in the sense used above: no numbers whose periods include all numbers smaller than the modulus and coprime to it. Yet an analogue is easy to see. A power of a number of the form $8k+3$ with odd exponent is always of the form $8k+3$, while one with even exponent is always of the form $8k+1$. So no power can have the form $8k+5$ or $8k+7$. Now the period of a number of the form $8k+3$ contains $2^{n-2}$ different terms, each of the form $8k+3$ or $8k+1$, and there are no more than $2^{n-2}$ such numbers smaller than the modulus. Thus every number of the form $8k+1$ or $8k+3$ is congruent modulo $2^n$ to some power of any number of the form $8k+3$. Similarly, the period of a number of the form $8k+5$ can be shown to contain all numbers of the forms $8k+1$ and $8k+5$. So if a number of the form $8k+5$ is taken as the base, all numbers of the forms $8k+1$ and $8k+5$, taken positively, and all numbers of the forms $8k+3$ and $8k+7$, taken negatively, acquire real indices. Indices congruent modulo $2^{n-2}$ must here be treated as equivalent. This is how to understand our Table I, where for the moduli 16, 32 and 64
\SourceLeaf{76}(no table is needed for the modulus $8$) we always took $5$ as the base. \emph{For example,} the number $19$, which is of the form $8n+3$ and must therefore be taken \emph{negatively}, has, for the modulus $64$, the index $7$. This means $5^7\equiv-19\,(\operatorname{mod.}64)$.

But numbers of the forms $8n+1$, $8n+5$, taken negatively, and those of the forms $8n+3$, $8n+7$, taken positively, would have to be assigned imaginary indices, so to speak. Introducing these reduces the calculation of indices to a very simple algorithm. But a fully rigorous treatment would take us too far afield, so we leave it for another occasion, when we may undertake a fuller account of the theory of imaginary quantities. In our view, no one has yet reduced that theory to clear concepts. Those with experience will readily find this algorithm for themselves. Those with less practice can still use the table just as people use logarithms without knowing recent authors’ discussions of imaginary \emph{logarithms}, provided they have a good grasp of the principles established above.

\begin{center}\small\emph{Moduli composed of several primes.}\end{center}

\Article{92}
For a modulus composed of several primes, almost everything about residues of powers can be derived from the general theory of congruences. But below we will show more fully how to reduce any congruence for such a modulus to congruences whose moduli are prime or powers of primes, so we need not spend much time on this here. We only note that the very fine property holding for the other moduli, namely that there always exist numbers whose periods include all numbers coprime to the modulus, fails here. There is just one exception: when the modulus is twice a prime or twice a power of a prime.
For if the modulus $m$ is put in the form $A^aB^bC^c$, etc., where $A,B,C$, etc., are distinct primes, and $A^{a-1}(A-1)$ is denoted by $\alpha$, $B^{b-1}(B-1)$ by $\beta$, etc., and finally $z$ is coprime to $m$, then $z^{\alpha}\equiv1\,(\operatorname{mod.}A^a)$. $z^{\beta}\equiv1\,(\operatorname{mod.}B^b)$, etc. So if $\mu$ is the least common multiple of $\alpha,\beta,\gamma$, etc., then $z^{\mu}\equiv1$ for all the moduli $A^a,B^b$, etc., and therefore also modulo $m$, their product. But except when $m$ is twice a prime or twice a power of a prime, the least common multiple of $\alpha,\beta,\gamma$, etc., is smaller than their product, since $\alpha,\beta,\gamma$, etc., cannot be pairwise coprime, but certainly have the common divisor $2$. So no number’s
\SourceLeaf{77}period can contain as many terms as there are numbers coprime to the modulus and smaller than it, since their number equals the product of $\alpha,\beta,\gamma$, etc. Thus, \emph{for example,} for $m=1001$ any number coprime to $m$, raised to the power with exponent $60$, is congruent to unity, since $60$ is a common multiple of $6,10,12$. --- But the case where the modulus is twice a prime or twice a power of a prime is entirely similar to the case of a prime or a power of a prime.

\Article{93}
We have already mentioned here and there the writings in which other mathematicians treated the subject of this section. But readers wanting some matters explained more fully than our wish for brevity allowed are referred particularly to the following papers by Euler. They are especially recommended for the clarity in which this great man always excelled above all others.

\emph{Theoremata circa residua ex divisione potestatum relicta Comm. nov. Petr. T.} VII, \emph{p.} 49 onward.

\emph{Demonstrationes circa residua ex divisione potestatum per numeros primos resultantia.}\par
\hspace*{\parindent}\emph{Ibid. T.} XVIII, \emph{p.} 85 onward.

To these may be added \emph{Opusculorum analyt. T.} I, papers 5 and 8.

\begin{center}\rule{.28\linewidth}{.25pt}\end{center}
